\slajdid{10-s032}
\begin{frame}[label=10-s032]
\gusttekst\setlength{\abovedisplayskip}{3pt}\setlength{\belowdisplayskip}{3pt}\setlength{\abovedisplayshortskip}{2pt}\setlength{\belowdisplayshortskip}{2pt}\begin{columns}[c,onlytextwidth]\column{.30\textwidth}\centering\begin{tikzpicture}[x=.65cm,y=.5cm,font=\sitneoznake,>=Latex]
\draw[->](0,4)--(5,4)node[below]{t};\draw[->](.2,3.7)--(.2,6.2)node[left]{$v_i(t)$};\draw(.2,5.7)--(1,5.7)--(1,4.3)--(2.8,4.3)--(2.8,5.7)--(4.7,5.7);\node[left]at(.2,5.7){$V_H$};\node[left]at(.2,4.3){$V_L$};
\draw[->](0,.7)--(5,.7)node[below]{t};\draw[->](.2,.4)--(.2,3.6)node[left]{$v_o(t)$};\foreach\y/\l in{2.8/V_{OH},1.9/50\%,1/V_{OL}}{\draw[dashed](.2,\y)--(3.3,\y);\node[left]at(.2,\y){$\l$};}
\draw(.2,1)--(1,1)..controls(1.2,2.5)and(1.7,2.8)..(2.2,2.8)--(2.8,2.8)..controls(2.85,1.8)and(2.95,1.2)..(3.05,1)--(4.7,1);
\draw[densely dotted](3.05,1)..controls(3.4,-1.3)and(3.6,-1.6)..(4.9,-1.6);\draw[->](1.6,-1.5)node[left]{po modelu}--(3.25,-.7);
\foreach\x/\l in{1/1,2.8/2}{\draw[dashed](\x,.45)--(\x,4.2);\node[below]at(\x,.45){$t_\l$};}
\draw[->,DEplava](1.9,-3.05)--(3.7,-1.7);
\end{tikzpicture}
\par\raggedright
Ne smemo koristiti gotove izraze za\par $t_{pHL}=0.69\tau$\par $t_f=2.2$
\column{.68\textwidth}
\[\begin{aligned}v_o(t_2+t_{pHL})&=V_{OH}-50\%(V_{OH}-V_{OL})\\&=V_{OH}-\frac{V_{OH}-V_{OL}}{2}=\frac{V_{OH}+V_{OL}}{2}\end{aligned}\]
pa je
\[t_{pHL}=\tau\ln\left(\frac{V_{OH}-(V_{CC}-\dfrac{R_C\beta_F}{R_B}(V_H-V_{BE}))}{\dfrac{V_{OH}+V_{OL}}2-(V_{CC}-\dfrac{R_C\beta_F}{R_B}(V_H-V_{BE}))}\right)\]
\[t_f=\tau\ln\left(\frac{(0.9V_{OH}+0.1V_{OL})-(V_{CC}-\dfrac{R_C\beta_F}{R_B}(V_H-V_{BE}))}{(0.1V_{OH}+0.9V_{OL})-(V_{CC}-\dfrac{R_C\beta_F}{R_B}(V_H-V_{BE}))}\right)\]
\end{columns}\medskip
\textcolor{DEcrvena}{Ostaje pitanje vrednosti $V_H$. Očigledno je da taj napon mora biti iz opsega $V_{IH}$ do $V_{OH}$, a na osnovu prethodnih diskusija kod strujnih kapaciteta izabraćemo najgori slučaj iz tog opsega, što pokazuju svi izvedeni izrazi, odnosno $V_H=V_{OHmin}$ a kao što je rečeno to bi i bilo u katalogu}
\end{frame}
